\section{Introduction}

% 背景
近年のAI技術の急速な発展は, 様々な分野において人間の能力を凌駕する成果を示しており, 人工知能が人間の知能を超えるという期待が社会全体で高まっている. 
GPT-4が米国司法試験で上位10\%のスコアを達成し \cite{openai2023gpt4}, 医療画像診断においてディープラーニングシステムが放射線科医の診断精度を上回り \cite{rajpurkar2017chexnet}, AlphaGoが囲碁の世界チャンピオンを破った \cite{silver2016mastering} など, 特定のタスクにおいてAIシステムが人間専門家の性能を凌駕する事例が相次いでいる. 
これらのAI技術の成功は, 主に大規模データセットの活用, 計算資源の増強, そしてアルゴリズムの継続的な改良によってもたらされている. 
このアプローチは工学的観点から極めて合理的であり, 実用的な成果を生み出し続けている. 

一方で興味深いことに, これほどAIが発展している現在においても, 人間の知能そのものについての理解は依然として限定的である. 
脳がどのように情報を処理し, 複雑な認知タスクを遂行しているのか, その根本的なメカニズムについては多くが未解明のまま残されている. 
例えば, 意識の神経基盤, binding problem, 量子マインド, 認知の具現化など, 多くの根本的な問題が未解明のまま残されている \cite{list2020unsolved}. 

% 問題提起/仮説
現在のAIシステムの性能評価は, 主として入力に対する出力の精度や速度といった表層的な指標に基づいて行われている. 
例えば機械翻訳では, BLEU \cite{papineni2002bleu} やMETEOR \cite{banerjee2005meteor} といった参照訳との表層的な一致度を測る指標が主流であり, 感情指向テキストにおける重要な翻訳エラーを検出する能力に限界がある \cite{saadany2021bleu}. 
このようなアプローチは, 特定のタスクの自動化や効率化という観点では有効であり, 産業応用においても大きな成果を上げていることは確かである. 

ここで重要な問いが生じる: もし我々が人間の脳の情報処理メカニズムをより深く理解し, それをAI技術に統合することができたならば, 現在のAIシステムをさらに進化させることができるのではないだろうか. 
人間の脳は, 現在の最先端AIシステムと比較して驚くべき効率性を示している. 
例えば, 人間の脳の消費電力は約20Wと推定されており \cite{clarke2008energy}, これは現代のGPUクラスターが消費する数キロワットと比較して桁違いに少ない. 
この驚異的なエネルギー効率で, 人間の脳は文脈理解, 創造的問題解決, 適応的学習といった高度な認知機能を実現している. 

真に人間社会に統合され, 人間と協調して働く実用的なAIシステムを構築するためには, 単なる入出力の最適化を超えた, より本質的な理解が必要となる. 
人間の脳がどのように情報を処理し, 文脈を理解し, 動的に戦略を選択しているのか, これらの原理を理解し取り入れることで, 単なる統計的パターンマッチングを超えた, より効率的で柔軟なAIシステムの実現が期待できる. 

% 研究の目的
本研究は, 同時通訳（simultaneous interpreting; SI）という, 人間と機械の双方にとって極めて挑戦的なタスクを研究対象として選択した. 
同時通訳が人間の認知メカニズム研究に最適である理由は複数存在する. 
第一に, 同時通訳は極めて高い速度と精度が同時に要求される高難易度の認知タスクである. 
通訳者は原発話の終了を待たずに対訳を産出する必要があり, 聴取, 理解, 翻訳, 発話という複数の認知プロセスを並列的に処理しなければならない. 
第二に, モダリティの観点から見ると, 同時通訳は単純な計算や情報分析とは異なり, 音声入力（聴取）, 意味理解, 短期・長期記憶との関連付け, 言語変換, 音声生成（発話）という複数のモダリティにまたがる統合的な認知処理を必要とする. 
第三に, 同時通訳は人間の専門家でも機械システムでも完全には達成できていない課題であり, 両者の処理戦略と限界を比較分析するのに理想的なタスクである. 

Gile \cite{gile1995regards} の努力モデルによれば, これらのプロセスは限られた認知資源を競合的に使用し, Seeber \cite{seeber2011cognitive} は特に構文的非対称性がある言語ペア（SVO--SOV）において認知負荷が著しく増加することを示している. 
Seeberの研究では, 同時通訳者が用いる4つの主要な認知戦略（waiting, stalling, chunking, anticipating）が特定されており, これらの戦略が認知負荷の管理において重要な役割を果たしていることが明らかになっている \cite{seeber2011cognitive}. 
脳神経科学の観点からは, Hervais-Adelman et al. \cite{hervais2017cortical} がfMRIを用いた研究により, 同時通訳時の脳活動パターンと神経ネットワークの特異性を明らかにしている. 

このように, 認知科学と脳神経科学の分野では, 人間の同時通訳者がどのような認知戦略を用い, どのような神経基盤に基づいて高度な通訳を実現しているかについて, 着実に理解が深まっている. 
しかしながら, これらの人間の処理メカニズムに関する貴重な知見は, 現在の同時通訳AIシステムの開発において十分に活用されているとは言い難い. 
その背景には, 前述の入力に対する出力の表層的な評価という問題が, 同時通訳の研究分野においても根強く存在していることがある. 

現在の同時通訳システムの評価においても, BLEUスコア \cite{papineni2002bleu} やSacreBLEU \cite{post2018call} といった表層的な一致度の評価に加え, 遅延時間（latency）と品質のトレードオフを考慮したSimulEval \cite{ma2020simuleval} や平均遅延（Average Lagging） \cite{ma2019stacl} など, 主に最終的な出力結果と時間的指標に基づく評価が中心となっている. 
これらの評価指標は確かに重要であるが, 人間の同時通訳者が実際にどのような認知処理を経て高品質な通訳を実現しているのか, その内部メカニズムについては考慮されていない. 
この評価パラダイムは, 現在の同時通訳AIシステムの開発においても, 単純なエンコーダ・デコーダモデルの最適化や, 待機時間の調整といった表層的な改善に留まる要因となっている. 
結果として, 認知科学や脳神経科学が明らかにしてきた人間の処理戦略の知見が, AI技術の発展に十分に反映されていないという状況が生じている. 

近年では, Meta's SeamlessM4T \cite{barrault2023seamlessm4t} やGoogle's Translatotron \cite{jia2019direct} など, エンドツーエンドの同時通訳システムが開発されている. 
SeamlessM4Tは初の統合型多言語・マルチモーダルAI翻訳モデルで, 最大100言語の音声・テキスト間の翻訳を単一モデルで実現するが, これらのシステムは依然として人間の同時通訳者が示す柔軟な戦略選択や文脈理解の深さには及ばない. 
人間の通訳者は通常15--30分で交代を必要とするという限界を持つ一方で, 状況に応じた適応的な処理戦略を駆使することで高品質な通訳を実現している. 

本研究では, 以下の仮説を提唱する: 人間の同時通訳研究を通じて明らかになっている認知戦略および神経処理メカニズムを, 最新のAI技術と体系的に比較分析することで, より発展的で効率的, かつ効果的なAIシステム開発のための重要な知見を得ることができる. 
具体的には, 人間の同時通訳者が示す適応的な処理戦略と, その背後にある神経基盤の理解を計算モデルに統合することで, 現在のAIシステムの限界を突破する可能性があると考える. 

本研究の目的は, 同時通訳タスクにおける人間側の処理メカニズムと機械側の処理メカニズムを横断的に比較分析することである. 
具体的には, 認知科学的観点から明らかにされている通訳者の認知戦略と, 脳神経科学的観点から解明されている神経基盤, そして情報工学的観点から開発されている同時通訳AIシステムの処理方式を体系的にレビューし, 比較検討を行う. 

この学際的な比較分析を通じて, 現在の同時通訳AIシステムをより理想的なものへと改善するための具体的な知見を導出することを目指す. 
本研究では, 人工知能の更なる発展において, 脳神経科学が明らかにする生理学的知見と, 認知科学が解明する人間の認知処理の原理を統合的に取り入れることで, 次世代の高度で実用的なAIシステムの実現が可能になるという立場を取る. 
人間の脳の処理戦略と機械の計算能力を組み合わせることで, 人間にとって自然でありながら, 現在の人間通訳者もAIシステムも実現できていない, 発展的な同時通訳システムの可能性を探求する. 